\scanpage{4}
\mspage{(2)}
que le $\mathbb Q$-type d'homotopie de $X$ est récupéré si $X$ est simplement connexe (\uline{sauf erreur}) -- de façon plus précise, il y a une équivalence de cat. \inserted{\unsure{donnée} par $C^{\unsure{*}}_{\mathrm{DRS}}$} entre la cat. hom. \inserted{faible} des espaces connexes et simplement connexes (du pt de vue singulier), et une certaine catégorie dérivée formée avec les $\mathbb Q$-algèbres graduées associatives et \inserted{anti} comm. \deleted{telles} à degrés $\geq 0$ telles que
\[
H^0(A)\simeq\mathbb Q\quad \deleted{\unsure{H^0(A)\simeq0}}\quad H^1(A)\simeq0.
\]
Il y a une théorie des modèles minimaux \inserted{pour de telles algèbres}, une façon très simple de récupérer les $\pi_i\otimes_{\mathbb Z}\mathbb Q$ en fonction d'un tel modèle... (On renvoie aux papiers de Sullivan)

Mais à nouveau, -- \unsure{perd} la torsion !

\figure{\begin{tikzpicture}\draw (0,0)--(14,0);\end{tikzpicture}}

Dém. du th. d'isomorphie de Sullivan
\[
H^q(C^\bullet_{\mathrm{DRS}\unsure{/\mathbb Q}}(X,\mathbb Q))\simeq H^q(X,\mathbb Q).
\]

$\forall n\geq0$, $[n]$ simplexe type de dim $n$

$\Delta^{[n]}_{\unsure{\mathbb R|*}}$ sa réalisation géométrique.
\margin{$(\sum_0^n X_i=1,\ \unsure{\forall}X_i\geq0)$}

$E^{[n]}$ l'espace affine qu'il sous-tend (avec une $\mathbb Q$-structure)
\margin{$\sum_0^n X_i=1$ dans $\mathbb R^{n+1}$}

$\unsure{\mathrm{DRS}^\bullet_{[n]}}=\unsure{C^\bullet_{\mathbb Q\mathrm{DR}}}(E^{[n]})$ son complexe de De Rh. -- $\mathbb Q$-algébrique.

C'est contravariant en $[n]$, d'où
\[
\mathrm{DRS}^\bullet_*=(\mathrm{DRS}_{[n]})_{n\geq0}
\]
Algèbre différentielle graduée \deleted{\unsure{semi}}\inserted{semi}-simpliciale (et $=$ simpliciale) -- à degrés $\geq0$, anticommutative.

Pour tt espace topologique $X$, $S_*(X)$ son complexe singulier $\frac12$ simplicial. On a
\[
C^\bullet_{\mathrm{DRS}}(X,\mathbb Q)\simeq\operatorname{Hom}(\unsure{S_*(X)},\mathrm{DRS}^\bullet_*)
\]
\mspage{3}
